\section{从订单明细到经营决策的任务链} % src:交接/计划.json:论文结构[文件名=论文/1.任务链与解题思路.tex].章节标题

\subsection{问题重述}

\textbf{问题1：}核算销售额分布与去重业务指标，刻画时间、地理、产品、客户和物流结构。 % src:交接/题面.json:问题[编号=1].要求

\textbf{问题2：}检验五类因素与销售额的稳定关联，区分增长驱动、经营短板及其解释边界。 % src:交接/题面.json:问题[编号=2].要求

\textbf{问题3：}用前三年训练月度预测模型，在后一年封存数据上检验误差与泛化能力。 % src:交接/题面.json:问题[编号=3].要求

\textbf{问题4：}预测下一年月、季和全年销售额，完成业务层级分解并形成经营动作。 % src:交接/题面.json:问题[编号=4].要求

\subsection{月度目标与双层证据链}

四问共用订单商品明细，却对应两种分析尺度：问题1、问题2保留商品、客户、区域和发货标签，用于解释销售结构；问题3、问题4按下单日期聚合自然月销售额，用于检验外推并配置经营总量。月度预测给出统一总量锚，明细层估计的稳定因素与收缩份额再把总量还原到各业务层级，因而任务链依次完成口径治理、驱动识别、总量预测和策略分解。 % src:交接/计划.json:问题清单[编号=1,2,3,4].输入,输出,依赖问题

我们最初按任务提示把序列视为约 $47$ 个月，重新解析日期并逐月聚合后确认研究期有 $48$ 个连续有单自然月，故采用 $36$ 个训练月和 $12$ 个测试月。 % src:交接/实验记录.json:[问题=总体规划/问题3].尝试,现象,决定
年度数据不足以训练并留出检验期，日度数据又会放大订单到达噪声；月度尺度既保留季节位置，也能直接汇总季度结果。 % src:交接/计划.json:问题清单[编号=3].方法理由

这条链路中的每次改造都来自实际复核：问题1曾把明细行直接视为订单，发现同单多行后改按订单号去重；问题2曾压缩为一单一行，但多品类订单会混合产品效应，遂保留明细并按订单成簇检验；问题3曾考虑自动定阶的大参数网格，短训练窗促使我们改用季节锚约束低自由度模型；问题4曾尝试为各细层单独建模，稀疏月与加总矛盾使方案转为总量锚定的份额收缩。 % src:交接/实验记录.json:[问题=问题1/订单口径],[问题=问题2/分析粒度与订单相关性],[问题=问题3/候选复杂度],[问题=问题4/层级预测方法选择]
